\documentclass[aspectratio=169,11pt]{beamer}

% --- Simple theme ---
\usetheme{default}
\usecolortheme{dove}
\setbeamertemplate{navigation symbols}{}
\setbeamertemplate{footline}[frame number]
\setbeamerfont{frametitle}{size=\large}

% --- Packages ---
\usepackage[utf8]{inputenc}
\usepackage[T1]{fontenc}
\usepackage{lmodern}
\usepackage{booktabs}
\usepackage{graphicx}
\usepackage{listings}
\usepackage{xcolor}
\usepackage{multirow}
\usepackage{textcomp}

% --- Code listing style ---
\definecolor{codebg}{gray}{0.96}
\definecolor{codegreen}{rgb}{0.2,0.5,0.2}
\definecolor{codegray}{rgb}{0.5,0.5,0.5}
\definecolor{codeblue}{rgb}{0.15,0.3,0.6}

\lstset{
  language=Python,
  basicstyle=\ttfamily\scriptsize,
  backgroundcolor=\color{codebg},
  keywordstyle=\color{codeblue}\bfseries,
  commentstyle=\color{codegreen},
  stringstyle=\color{codegray},
  showstringspaces=false,
  breaklines=true,
  frame=single,
  framerule=0pt,
  xleftmargin=4pt,
  xrightmargin=4pt,
  aboveskip=6pt,
  belowskip=4pt,
}

% --- Title ---
\title{SOM-GNN-SDM Switzerland}
\subtitle{Species Distribution Modelling with Graph Neural Networks\\at 30\,m Resolution}
\author{Marco Kronenberg}
\date{CAS AML --- May 2026}

\begin{document}

% ============================================================
\begin{frame}
\titlepage
\end{frame}

% ============================================================
\begin{frame}{Motivation --- Why Graphs for SDM?}

\textbf{Problem:} Traditional SDMs (MaxEnt, RF) treat each location independently.\\
No spatial context --- a pixel's prediction ignores its neighbours.

\bigskip

\textbf{Insight:} Habitat suitability is \emph{not} independent.\\
A patch surrounded by suitable habitat is more valuable (connectivity, dispersal).

\bigskip

\textbf{GNN-SDM} (Wu et al.\ 2025): Encode the landscape as a graph.
\begin{itemize}
  \item Nodes = habitat patches (with environmental features)
  \item Edges = spatial adjacency
  \item Message passing lets each node ``see'' its neighbourhood
\end{itemize}

\bigskip

\textbf{This project:} Adapt the framework to Swiss flora at 30\,m resolution.
\begin{itemize}
  \item Wu et al.: uniform 10\textquotesingle{} grid ($\sim$18.5\,km cells), continental scale
  \item Ours: SOM-derived patches from 30\,m pixels
  \item 3,756 plant species from GBIF, 21 environmental features
\end{itemize}

\end{frame}

% ============================================================
\begin{frame}{Pipeline --- From Pixels to Graph}

\begin{center}
\small
\begin{tabular}{ccccc}
\fbox{\parbox{2.2cm}{\centering 300M pixels\\21 features}}
& $\longrightarrow$ &
\fbox{\parbox{2.4cm}{\centering SOM (64$\times$64)\\$\rightarrow$ 30 types}}
& $\longrightarrow$ &
\fbox{\parbox{2.8cm}{\centering Connected comp.\\166,464 patches}}
\\[6pt]
& & & $\longrightarrow$ &
\fbox{\parbox{2.4cm}{\centering Graph\\465K edges}}
\end{tabular}
\end{center}

\bigskip

\textbf{12 features per node} (selected from 21 via RF importance):

\smallskip

{\small
\begin{tabular}{llll}
\toprule
\textbf{Group} & \textbf{Features} & \textbf{Source} & \textbf{Native res.} \\
\midrule
Terrain & elevation, slope, TWI & Copernicus GLO-30 & 30\,m \\
Climate & bio1, bio4, bio12 & CHELSAch 1981--2010 & 25\,m \\
Land cover & tree, grass, built-up fractions & ESA WorldCover 2021 & 10\,m \\
Biological & canopy height, NDVI, HFP & ETH / GIMMS / Mu et al. & 10\,m--1\,km \\
\bottomrule
\end{tabular}
}

\smallskip

All regridded to 30\,m master grid, then mean-aggregated per patch $\rightarrow$ 12-dim node vector.

\end{frame}

% ============================================================
\begin{frame}{Feature Correlation Matrix}

\begin{center}
\includegraphics[height=0.85\textheight]{figures/correlation_matrix.png}
\end{center}

\end{frame}

% ============================================================
\begin{frame}{Landscape Patches --- SOM-Based Segmentation}

\begin{center}
\includegraphics[height=0.85\textheight]{figures/patch_map.png}
\end{center}

\end{frame}

% ============================================================
\begin{frame}[fragile]{Model --- GraphSAGE Architecture}

\begin{lstlisting}
class GNNSDM(torch.nn.Module):
    def __init__(self, in_channels, hidden_dims=[64, 48, 32], dropout=0.2):
        super().__init__()
        self.convs = torch.nn.ModuleList()
        prev = in_channels
        for dim in hidden_dims:
            self.convs.append(SAGEConv(prev, dim, aggr='mean'))
            prev = dim
        self.out = torch.nn.Linear(prev, 1)

    def forward(self, x, edge_index):
        for conv in self.convs:
            x = F.leaky_relu(conv(x, edge_index))
            x = F.dropout(x, p=self.dropout, training=self.training)
        return torch.sigmoid(self.out(x)).squeeze(-1)
\end{lstlisting}

\textbf{Key design choices:}
\begin{itemize}\setlength\itemsep{1pt}
  \item Architecture search on 6 species $\rightarrow$ \texttt{[64, 48, 32]} beats paper's \texttt{[24, 18, 8]}
  \item Random pseudo-absences (3:1 ratio) --- same as RF for fair comparison
  \item Node weighting via PageRank; weighted MSE loss; early stopping (patience=50)
\end{itemize}

\end{frame}

% ============================================================
\begin{frame}{Results --- 12 Species Deep Dive (GNN vs RF)}

{\small
\begin{tabular}{llrrrc}
\toprule
\textbf{Species} & \textbf{Group} & \textbf{Patches} & \textbf{GNN AUC} & \textbf{RF AUC} & \textbf{$\Delta$} \\
\midrule
\textit{Senecio inaequidens}  & common     & 4,838  & \textbf{0.917} & 0.910 & +0.007 \\
\textit{Fagus sylvatica}      & common     & 14,305 & \textbf{0.822} & 0.788 & +0.034 \\
\textit{Picea abies}          & common     & 7,526  & \textbf{0.801} & 0.773 & +0.028 \\
\textit{Phragmites australis} & common     & 3,662  & \textbf{0.800} & 0.787 & +0.013 \\
\textit{Potentilla erecta}    & common     & 5,527  & \textbf{0.774} & 0.760 & +0.014 \\
\textit{Trifolium pratense}   & common     & 7,383  & \textbf{0.757} & 0.720 & +0.037 \\
\midrule
\textit{Aquilegia alpina}     & vulnerable & 679    & \textbf{0.906} & 0.879 & +0.027 \\
\textit{Pulsatilla vulgaris}  & vulnerable & 550    & 0.844 & 0.851 & $-$0.007 \\
\textit{Gladiolus palustris}  & vulnerable & 329    & 0.838 & 0.838 & $\approx$0 \\
\textit{Stipa pennata}        & vulnerable & 392    & 0.834 & 0.844 & $-$0.010 \\
\textit{Cypripedium calceolus}& vulnerable & 2,423  & \textbf{0.818} & 0.797 & +0.021 \\
\textit{Drosera rotundifolia} & vulnerable & 1,499  & 0.812 & 0.812 & $\approx$0 \\
\bottomrule
\end{tabular}
}

\bigskip

\begin{itemize}\setlength\itemsep{1pt}
  \item GNN consistently outperforms RF on common species (mean $\Delta$ = +0.022)
  \item Vulnerable species: mixed --- GNN wins where more presence patches exist
\end{itemize}

\end{frame}

% ============================================================
\begin{frame}{Predicted vs Observed Prevalence}

\begin{center}
\includegraphics[width=\textwidth,height=0.82\textheight,keepaspectratio]{figures/prevalence.png}
\end{center}

\end{frame}

% ============================================================
\begin{frame}[fragile]{Next Steps --- Conservation Application}

\textbf{What's done:}
\begin{itemize}\setlength\itemsep{1pt}
  \item[\checkmark] Full pipeline: data $\rightarrow$ features $\rightarrow$ SOM $\rightarrow$ patches $\rightarrow$ graph $\rightarrow$ GNN
  \item[\checkmark] 3,751 species with suitability scores
  \item[\checkmark] RF baseline for comparison (same patches, no graph)
  \item[\checkmark] IUCN threat status for all species
\end{itemize}

\medskip

\textbf{What's next:}
\begin{enumerate}\setlength\itemsep{1pt}
  \item \textbf{GNN vs RF full comparison} --- quantify where graph structure helps
  \item \textbf{Circuit-theory connectivity} --- current flow through suitable habitat
  \item \textbf{Gap analysis} --- overlay with protected areas, find coverage gaps
\end{enumerate}

\begin{lstlisting}
# Circuit theory: resistance = 1 - suitability
# -> current flow identifies critical corridors
priority_score = (mean_suitability + mean_current_flow) / 2
\end{lstlisting}

\textbf{Open questions:}
\begin{itemize}\setlength\itemsep{1pt}
  \item Does GNN help more for rare species? (connectivity matters more?)
  \item Can we lower the 100-record threshold for threatened species?
\end{itemize}

\end{frame}

% ============================================================
\end{document}
